\documentclass[11pt]{article}
\usepackage{amsmath,amssymb,amsthm,mathtools,geometry,hyperref}
\geometry{margin=1in}
\newcommand{\C}{\mathbb C}
\newcommand{\Res}{\operatorname{Res}}
\newcommand{\diag}{\operatorname{diag}}
\title{Arithmetic residues on the complete ES--Fable factor space}
\author{Independent programme calculation}
\date{20 September 2026}
\begin{document}
\maketitle
This note independently checks equations (1)--(17) of the incoming
calculation \cite{Incoming}. It retains the original quartic
coefficient $-1/2$, both factor signs, the projective root at
infinity and the original polynomial coordinates. It corrects
the incoming unnumbered zeta determinant following (16).
No numerical zero calculation is used.

\section{The complete original factors and arithmetic residue map}
In original coordinates $q=(a,y,z,w)$ the polynomial is
\begin{align}
b&=i+ay,&c&=-i+2ay+a^2z,\nonumber\\
d&=-iy-a(iz+2y^2)-a^2yz,&e&=2z-7iy^2+aw,\nonumber\\
f&=iw+3iy^3-4yz+a(6iy^2z+wy+4y^4)+2a^2y^3z,\nonumber\\
P(q)&=(ac,ae+bd,af+be,bf). \tag{AR1}
\end{align}
Substitution proves $ad+bc=1$ and $C(-b,a)=1$, where
$L=aU+bV$ and $C=cU^3+dU^2V+eUV^2+fV^3$.
The full factor space and its exact marked-root morphism are
\begin{align}
{\cal E}&=\{ad+bc=1,\ a^3f-a^2be+ab^2d-b^3c=1\},\nonumber\\
{\cal P}&=(ac,ae+bd,af+be,bf),\qquad
{\cal E}\longrightarrow\{(u,[U:V]):H_u(U,V)=0\},\nonumber\\
(a,b,c,d,e,f)&\longmapsto({\cal P},[-b:a]),\qquad
H_u=AU^4+U^3V+BU^2V^2+C_1UV^3+DV^4.
\tag{AR2}
\end{align}
The first equation ensures $(a,b)\ne(0,0)$. Simultaneously
negating all six factors leaves the target and marked root fixed.
For a finite simple root $r$ of
$h_u(s)=As^4+s^3+Bs^2+C_1s+D$, division by $a(U-rV)$ gives
\[
b=-ar,\quad c=A/a,\quad d=(1+Ar)/a,\quad
e=(B+r+Ar^2)/a,\quad f=(C_1+Br+r^2+Ar^3)/a.
\]
The resultant becomes $a^2h'_u(r)=1$. This proves that both
square roots give exactly the two full factor states over a
simple finite root, while a multiple root gives none.
For the original completed function \cite{DLMFXi},
\begin{equation}
r=-y-i/a,\quad
d(\xi\circ r)=\xi'(r)(ia^{-2}da-dy),\qquad
\Res_r\frac{\xi(s)\,ds}{h_u(s)}
=\frac{\xi(r)}{h'_u(r)}=a^2\xi(r).
\tag{AR3}
\end{equation}
Thus each pair of factor signs gives one arithmetic value
and one residue, with no sign-dependent zero test.

For the other projective chart, put $t=V/U$, $g_u(t)=H_u(1,t)$
and $t_0=-a/b$. Division by $b(V-t_0U)$ gives
\begin{equation}
-b^2g'_u(t_0)=1,\qquad
\frac{\xi(s)\,ds}{h_u(s)}
=-\frac{t^2\xi(1/t)}{g_u(t)}\,dt.
\tag{AR4}
\end{equation}
Indeed $C(-b,a)=(-b)^3C(1,t_0)=-b^2g'_u(t_0)$.
For $t_0\ne0$ the residue is
$b^2t_0^2\xi(1/t_0)=a^2\xi(r)$, proving exact equality
of the chart formulas. At $A=0$ the root $t_0=0$ is simple
because $g'_u(0)=1$. Its two full factor states are
\begin{equation}
(a,b,c,d,e,f)=
(0,\sigma i,-\sigma i,-\sigma iB,-\sigma iC_1,-\sigma iD),
\qquad \sigma=\pm1.
\tag{AR5}
\end{equation}
Only $\sigma=1$ occurs in AR1, since $a=0$ forces $b=i$.
Its original coordinates are
$(0,B,i(7B^2-C_1)/2,11B^3-2BC_1-D)$.

\section{The exact finite remainder and infinity residue}
Let $n=\deg h_u$, equal to four or three, and let $\alpha_n$
be its actual leading coefficient, respectively $A$ or $1$.
For each entire function $F$ there is exactly one polynomial
$R_F$ of degree less than $n$ such that
\begin{equation}
F=R_F+h_uG_F,\quad G_F\text{ entire},\qquad
R_F^{(j)}(r)=F^{(j)}(r)\quad(0\le j<m_r),\qquad
{\cal A}_u=\C[s]/(h_u).
\tag{AR6}
\end{equation}
Here $m_r$ is the root multiplicity. For a complete existence
proof, the ideals $(s-r)^{m_r}$ are pairwise coprime; the
Euclidean algorithm gives their Bezout identities and the
polynomial Chinese remainder construction. Prescribing the
displayed Taylor data and reducing modulo $h_u$ gives a
representative of degree less than $n$. Its quotient in AR6
has only removable singularities. Uniqueness follows because
a polynomial of degree less than $n$ divisible by $h_u$ is
zero. Hence $j_uF=[R_F]$ is an algebra morphism.
In each local Taylor basis its multiplication matrix is
triangular with $F(r)$ on the diagonal $m_r$ times, so
\begin{equation}
\det M_{j_uF}=\prod_rF(r)^{m_r},\qquad
j_uF\text{ is a unit exactly when every }F(r)\ne0.
\tag{AR7}
\end{equation}
This includes nonreduced root schemes, without asserting
existence of a finite factor state at their multiple roots.

Write $h_u=(s-r)^{m_r}k_r(s)$, $k_r(r)\ne0$. The full
principal part and the infinity residue are
\begin{equation}
\left[\frac{F}{h_u}\right]_{\mathrm{pp},r}
=\sum_{j=1}^{m_r}\frac{(F/k_r)^{(m_r-j)}(r)}{(m_r-j)!}
 (s-r)^{-j},\qquad
\Res_\infty\frac{F\,ds}{h_u}
=-\frac{[s^{n-1}]R_F}{\alpha_n}.
\tag{AR8}
\end{equation}
The first formula is Taylor expansion. For the second,
$F/h_u=R_F/h_u+G_F$; an entire $G_F$ contributes no
$s^{-1}$ coefficient at infinity. The rational first term
has coefficient $[s^{n-1}]R_F/\alpha_n$, and $t=1/s$
gives the minus sign. Equivalently this residue is minus
the sum of all finite residues. It has the same value at
both infinity signs AR5.

There is no holomorphic point value $\xi(\infty)$.
Indeed $\xi(2m)=m(2m-1)\pi^{-m}(m-1)!\zeta(2m)$ and
$\zeta(2m)>1$ by its positive Dirichlet series. The factorial
grows faster than any fixed power, so $\xi$ is not a polynomial.
An entire function meromorphic at infinity would have a
polynomial growth bound, and Cauchy's coefficient formula would
make it a polynomial. Thus $\xi(1/t)$ is essential at $0$.
AR8 supplies a differential residue, not an invented value
or local zero test at infinity.

At $h_0=s^3-s$, the completion and reflection give
$\xi(0)=\xi(1)=1/2$ and $\xi(-1)=\xi(2)=\pi/6$.
Evaluation at the three roots proves
\begin{align}
R_\xi&=\tfrac12+\frac{\pi-3}{12}(s^2-s),&
R_\xi^{-1}&=2+(3/\pi-1)(s^2-s)\pmod{s^3-s},\nonumber\\
\det M_{R_\xi}&=\pi/24,&
(\Res_0,\Res_1,\Res_{-1},\Res_\infty)\frac{\xi\,ds}{s^3-s}
&=(-\tfrac12,\tfrac14,\tfrac{\pi}{12},-\tfrac{\pi-3}{12}).
\tag{AR9}
\end{align}
Each value occurs at both full factor signs; only one of the
infinity signs is in the original chart. All are nonzero.

\section{The literal quartet and every completion factor}
Fix $0<\delta<1/2$, $\gamma>0$, and retain
$r_{\epsilon,\eta}=1/2+\epsilon\delta+i\eta\gamma$,
with original order $(+,+),(+,-),(-,+),(-,-)$.
Put $\rho=r_{+,+}$, $x=s-1/2$, $k=\delta^2-\gamma^2$,
$l=2\delta\gamma$, $H_\delta=(x^2-k)^2+l^2$ and
$h_\rho=-H_\delta/2$. The last factor is forced by the
original cubic coefficient one. Write
$Z=\xi(\rho)=U+iV$, $\alpha=V/l$, $\beta=U-k\alpha$.
Reflection and conjugation give the four values
$(Z,\overline Z,\overline Z,Z)$, so
\begin{equation}
R_\xi=\alpha x^2+\beta=U+\alpha(x^2-k),\qquad
\det M_{R_\xi}=|Z|^4,\qquad
R_\xi^{-1}=\frac{U-\alpha(x^2-k)}{|Z|^2}\quad(Z\ne0).
\tag{AR10}
\end{equation}
Evaluation proves the remainder by AR6. The identity
$(x^2-k)^2=-l^2$ in the quotient proves the inverse.
For both original sign states at every root,
\begin{equation}
h'_\rho(r_{\epsilon,\eta})
=4\delta\gamma(\epsilon\gamma-i\eta\delta),\quad
a_{\epsilon,\eta,\pm}^2
=\frac{1}{4\delta\gamma(\epsilon\gamma-i\eta\delta)},\quad
\Res_r\frac{\xi\,ds}{h_\rho}=a^2\xi(r),\quad
\Res_r\frac{2\xi\,ds}{H_\delta}=-a^2\xi(r).
\tag{AR11}
\end{equation}
Multiplication of the three root differences and $-1/2$
proves the derivative, retaining the original factor and sign.

On the open critical strip the exact completion is
$\xi(s)=\Phi(s)\zeta(s)$ with
$\Phi(s)=s(s-1)\pi^{-s/2}\Gamma(s/2)/2$.
It is holomorphic and nonzero there: the polynomial factors
do not vanish, and Gamma has no poles or zeros in the right
half-plane. Thus its jet is a unit, including at the doubled
quartet below, and the complete multiplication morphism is
\begin{equation}
M_{j_u\xi}=M_{j_u\Phi}M_{j_u\zeta},\qquad
M_{j_u\zeta}=M_{j_u(1/\Phi)}M_{j_u\xi}.
\tag{AR12}
\end{equation}
This is an invertible map in the same coefficient algebra,
not a replacement of the original coordinates or norm.
It is not a global unit assertion: $\Phi(0)=-1$,
$\Phi(1)=0$, $\Phi'(1)=1/2$, and $\Phi$ has poles at the
negative even integers. The zeta pole and trivial zeros
compensate these factors in the entire completion.
AR12 therefore is not applied to the root $1$ of AR9 or to
the projective root at infinity.

The exact zeta determinant is
\begin{align}
\det M_{j_{H_\delta}\zeta}
&=|\zeta(\rho)\zeta(1-\rho)|^2
=\frac{|Z|^4}{|\Phi(\rho)\Phi(1-\rho)|^2}\nonumber\\
&=\left|\frac{\Phi(\rho)}{\Phi(1-\rho)}\right|^2
                         |\zeta(\rho)|^4,\nonumber\\
\frac{\Phi(\rho)}{\Phi(1-\rho)}
&=\pi^{1/2-\rho}\frac{\Gamma(\rho/2)}{\Gamma((1-\rho)/2)},\nonumber\\
\left|\frac{\Phi(\rho)}{\Phi(1-\rho)}\right|^2
&=\pi^{-2\delta}
\frac{|\Gamma(1/4+\delta/2+i\gamma/2)|^2}
     {|\Gamma(1/4-\delta/2+i\gamma/2)|^2}.
\tag{AR13}
\end{align}
Indeed the four values are
$\zeta(\rho),\overline{\zeta(\rho)},
\overline{\zeta(1-\rho)},\zeta(1-\rho)$.
Apply AR7 and then $\xi(\rho)=\xi(1-\rho)$, dividing only
by the nonzero completion factors, never by a possible zero
of zeta. The polynomial factors cancel since
$(1-s)(-s)=s(s-1)$, leaving the displayed Gamma and $\pi$
factors. The incoming $|\zeta(\rho)|^4$ formula omits this
generally nonunit-modulus ratio. It is not an identity:
as $\rho\to1$ near the real axis inside the strip, the ratio
tends to zero while $\xi(\rho)\to1/2$. Sufficiently nearby
nonreal points have nonzero zeta values and ratio modulus
different from one.

\section{Confluence in the unchanged coefficient coordinates}
Write $H_\delta=s^4+c_3s^3+c_2s^2+c_1s+c_0$, where
\[
c_3=-2,\quad c_2=\tfrac32+2(\gamma^2-\delta^2),\quad
c_1=-\tfrac12-2(\gamma^2-\delta^2),\quad
c_0=\tfrac1{16}+\tfrac{\gamma^2-\delta^2}{2}
                              +(\delta^2+\gamma^2)^2.
\]
In the original basis $(1,s,s^2,s^3)$, multiplication by
$s$ and by the completed remainder are exactly
\begin{equation}
C_\delta=\begin{pmatrix}
0&0&0&-c_0\\1&0&0&-c_1\\0&1&0&-c_2\\0&0&1&2
\end{pmatrix},\qquad
[M_{j\xi}]=\alpha(C_\delta-\tfrac12I)^2+\beta I .
\tag{AR14}
\end{equation}
The last column is the defining polynomial relation; the
others multiply the first three original monomials.
Let $r_\eta=1/2+i\eta\gamma$. Retain the exact paired
evaluation morphism
\[
{\cal V}_\delta p=
\left(\frac{p(r_++\delta)+p(r_+-\delta)}2,\
\frac{p(r_++\delta)-p(r_+-\delta)}{2\delta},\
\frac{p(r_-+\delta)+p(r_--\delta)}2,\
\frac{p(r_-+\delta)-p(r_--\delta)}{2\delta}\right)^{\mathsf T}.
\]
Its original-coordinate rows and determinant are
\begin{equation}
(1,r_\eta,r_\eta^2+\delta^2,r_\eta^3+3r_\eta\delta^2),\quad
(0,1,2r_\eta,3r_\eta^2+\delta^2),\qquad
\det{\cal V}_\delta=16\gamma^2(\gamma^2+\delta^2).
\tag{AR15}
\end{equation}
Each pair recovers both original evaluations by $u\pm\delta v$.
The complete inverse in original coefficient coordinates, for
$\delta>0$, is the polynomial
\[
{\cal V}_\delta^{-1}(u_+,v_+,u_-,v_-)(s)
=\sum_{\eta=\pm1}\sum_{\epsilon=\pm1}
(u_\eta+\epsilon\delta v_\eta)
\frac{H_\delta(s)}
 {(s-r_{\epsilon,\eta})H'_\delta(r_{\epsilon,\eta})}.
\]
Every quotient in this sum is a polynomial of degree three;
evaluation at each marked root proves both inverse compositions.
Its product law is $(u,v)(u',v')=(uu'+\delta^2vv',uv'+vu')$,
the algebra $\C[t]/(t^2-\delta^2)$. The Vandermonde and the
two explicit pair changes give the determinant; direct
expansion gives the same expression. It remains nonzero
at $\delta=0$. No root value or factor sign has been discarded.
At zero the map is
$ {\cal V}_0p=(p(r_+),p'(r_+),p(r_-),p'(r_-))$.
Put $d_\eta=r_\eta-r_{-\eta}$ and
\[
L_\eta(s)=(s-r_{-\eta})/d_\eta,\quad
A_\eta(s)=(1-2(s-r_\eta)/d_\eta)L_\eta(s)^2,\quad
B_\eta(s)=(s-r_\eta)L_\eta(s)^2.
\]
Direct evaluation and differentiation prove its full inverse
\begin{equation}
{\cal V}_0^{-1}(u_+,v_+,u_-,v_-)
=u_+A_++v_+B_++u_-A_-+v_-B_-.
\tag{AR16}
\end{equation}
For any holomorphic $F$, its local multiplication block is
$\begin{psmallmatrix}F(r_\eta)&0\\F'(r_\eta)&F(r_\eta)\end{psmallmatrix}$.
Conjugation by the explicit $ {\cal V}_0$ gives the matrix
back in the original coefficient basis.
Put $X(\gamma)=\xi(1/2+i\gamma)\in\mathbb R$.
Differentiation gives $\xi'(r_\eta)=-i\eta X'$;
Taylor expansion in real $\delta$ gives $U\to X$,
$V/\delta\to-X'$. Hence the full limiting data are
\begin{align}
H_0&=((s-\tfrac12)^2+\gamma^2)^2,\quad
T=(s-\tfrac12)^2+\gamma^2,\quad T^2=0\text{ in }{\cal A}_0,\nonumber\\
R_{\xi,0}&=X-\frac{X'}{2\gamma}T,\qquad
{\cal V}_0M_{j\xi}{\cal V}_0^{-1}
=\diag\left(\begin{pmatrix}X&0\\-iX'&X\end{pmatrix},
            \begin{pmatrix}X&0\\iX'&X\end{pmatrix}\right),\nonumber\\
\det M_{j\xi}&=X^4,\qquad
R_{\xi,0}^{-1}=X^{-1}+\frac{X'}{2\gamma X^2}T\quad(X\ne0).
\tag{AR17}
\end{align}
The same identities follow from AR16 using
$T(r_\eta)=0$, $T'(r_\eta)=2i\eta\gamma$.
If $X'\ne0$, there are two Jordan blocks of size two with
eigenvalue $X$. If $X'=0$, the operator is exactly $XI$.
At $X=0$, $X'\ne0$, its rank is two and its square is zero;
at $X=X'=0$, it is zero. These exhaustive alternatives follow
from the matrices, without any assumed zero multiplicity.
Thus $H_0$ divides $\xi$ exactly when $X=X'=0$.

\section{Completion jets and the converging residues}
Let $\phi_\eta=\Phi(r_\eta)\ne0$, $z_\eta=\zeta(r_\eta)$.
The complete original completion data are
\begin{align}
L_\Phi(s)&=\frac{\Phi'(s)}{\Phi(s)}
=\frac1s+\frac1{s-1}-\frac{\log\pi}{2}
                  +\frac12\frac{\Gamma'(s/2)}{\Gamma(s/2)},\nonumber\\
{\cal V}_0M_\Phi{\cal V}_0^{-1}\big|_\eta
&=\begin{pmatrix}\phi_\eta&0\\\phi_\eta L_\Phi(r_\eta)&\phi_\eta\end{pmatrix},
\quad
({\cal V}_0M_\Phi{\cal V}_0^{-1})^{-1}\big|_\eta
=\begin{pmatrix}\phi_\eta^{-1}&0\\
-L_\Phi(r_\eta)\phi_\eta^{-1}&\phi_\eta^{-1}\end{pmatrix},\nonumber\\
z_\eta&=X/\phi_\eta,\quad
z_\eta'=\frac{-i\eta X'-XL_\Phi(r_\eta)}{\phi_\eta},\qquad
\det M_{j_{H_0}\zeta}=|z_+|^4=\frac{X^4}{|\phi_+|^4}.
\tag{AX1}
\end{align}
Differentiate the original product defining $\Phi$ to prove its
logarithmic derivative. Multiplication of the triangular matrices
proves their inverse. Conjugation gives
$\phi_-=\overline{\phi_+}$ and the determinant. Thus the bare
$|\zeta(\rho)|^4$ formula holds at the doubled critical-line
limit, but needs AR13 at the distinct off-line quartet.
The exact unit map preserves the vanishing order even at $X=0$.

At $r_\eta$, the original polynomial $h_{\rho,0}=-H_0/2$
has complete arithmetic principal part
\begin{equation}
\left[\frac{\xi(s)}{h_{\rho,0}(s)}\right]_{\mathrm{pp},r_\eta}
=\frac{X}{2\gamma^2}(s-r_\eta)^{-2}
+\frac{i\eta(X-\gamma X')}{2\gamma^3}(s-r_\eta)^{-1}.
\tag{AX2}
\end{equation}
Expand $-2\xi(s)/(s-r_{-\eta})^2$ at $r_\eta$.
Its constant and linear terms are
$-2X/d_\eta^2$ and
$-2\xi'(r_\eta)/d_\eta^2+4X/d_\eta^3$.
Substitution $d_\eta=2i\eta\gamma$ proves every sign in AX2.

For $\delta>0$ denote the two finite-root residues at fixed
$\eta$ by $ {\cal R}_{\epsilon,\eta}$; each still occurs twice
among the signed states. AR11 and Taylor expansion give
\begin{align}
{\cal R}_{\epsilon,\eta}(\delta)
&=\frac{\epsilon X}{4\delta\gamma^2}
  +\frac{i\eta(X-\gamma X')}{4\gamma^3}+O(\delta),\nonumber\\
\sum_{\epsilon=\pm1}{\cal R}_{\epsilon,\eta}(\delta)
&\longrightarrow\frac{i\eta(X-\gamma X')}{2\gamma^3},&
\sum_{\epsilon=\pm1}\epsilon\delta\,{\cal R}_{\epsilon,\eta}(\delta)
&\longrightarrow\frac{X}{2\gamma^2}.
\tag{AX3}
\end{align}
The remainder is a convergent local Taylor remainder for fixed
$\gamma>0$. The exact rational identity
\[
\sum_{\epsilon=\pm1}\frac{{\cal R}_{\epsilon,\eta}}
 {s-r_\eta-\epsilon\delta}
=\frac{(s-r_\eta)({\cal R}_{+,\eta}+{\cal R}_{-,\eta})
       +\delta({\cal R}_{+,\eta}-{\cal R}_{-,\eta})}
       {(s-r_\eta)^2-\delta^2}
\]
then gives precisely AX2, locally uniformly away from $r_\eta$.
At $X=0$, $X'\ne0$, the individual residues tend to
$-i\eta X'/(4\gamma^2)$ while their factors $a$ diverge.
No finite full factor state exists over the doubled roots:
$a^2h'(r)=1$ excludes every finite root and $A=-1/2$ excludes
an infinity root. AR6, AR16 and AX2 nevertheless give the exact
jet and residue limit of these escaping factors.

\section{The whole explicit escape family: an exact arithmetic exclusion}
For every real $\varepsilon>0$, the original source
$(\varepsilon,0,i/2,0)$ has the exact target polynomial
\begin{equation}
h_\varepsilon(s)=i\varepsilon s(s+i/\varepsilon)
 [(\varepsilon^2/2-1)s^2-i\varepsilon s/2+1],\qquad
r_\pm=\frac{i\varepsilon\pm\sqrt{16-9\varepsilon^2}}
                 {2(\varepsilon^2-2)}.
\tag{AX4}
\end{equation}
Multiplication gives target
$(-i\varepsilon+i\varepsilon^3/2,3i\varepsilon/2,-1,0)$.
The other two labels are $0,-i/\varepsilon$.
At $\varepsilon^2=2$ the last factor is linear, its root is
$-2i/\varepsilon$, and the fourth projective root is infinity.
At $\varepsilon^2=16/9$ the two quadratic roots coincide
and are purely imaginary; for larger squares they are also
purely imaginary. When $0<\varepsilon^2<16/9$, their real
parts are opposite. The positive part equals
$R=\sqrt{16-9\varepsilon^2}/[2(2-\varepsilon^2)]$.
Squaring positive denominators proves $0<R<1$ exactly when
$7/4<\varepsilon^2<16/9$. There
$|\Im r_\pm|=\varepsilon/[2(2-\varepsilon^2)]<3$.
This covers all parameters, including changes of degree
and multiplicity.

Here is the complete classical boundary argument used for
these labels. For $\sigma>1$, the absolutely convergent Euler
logarithm gives, when $t\ne0$,
\[
3\log|\zeta(\sigma)|+4\log|\zeta(\sigma+it)|
 +\log|\zeta(\sigma+2it)|
=\sum_{p,m\ge1}
\frac{3+4\cos(mt\log p)+\cos(2mt\log p)}{mp^{m\sigma}}\ge0,
\]
since the numerator is $2(1+\cos(mt\log p))^2$.
If $\zeta$ had a zero of order $m_0\ge1$ at $1+it$,
the product obtained by exponentiating would be
$O((\sigma-1)^{4m_0-3})\to0$, using the simple pole
at $1$ and boundedness at $1+2it$. This contradicts its
lower bound one. Thus zeta has no zeros on $\Re s=1$;
the Euler product gives nonvanishing for $\Re s>1$.
The original completion factors are nonzero there away
from $1$, where $\xi(1)=1/2$. Reflection gives the same
for $\xi$ on $\Re s\le0$. This is the classical
Hadamard--de la Vallée Poussin line-one argument, reproduced
here and not claimed as a new theorem.

The precise theta input is the original series and modular
transformation in \cite{DLMFTheta}. At $z=0,\tau=ix$ these give,
with $\psi(x)=\sum_{n\ge1}e^{-\pi n^2x}$,
$1+2\psi(x)=x^{-1/2}(1+2\psi(1/x))$.
For $\Re s>1$, termwise integration gives
$\int_0^\infty\psi(x)x^{s/2-1}dx
=\Gamma(s/2)\pi^{-s/2}\zeta(s)$.
Split the integral at one and substitute $x=1/y$ in its
lower part. The elementary contribution is
$1/(s-1)-1/s=1/[s(s-1)]$.
Multiplying by $s(s-1)/2$ and using the completion gives
\[
\xi(s)=\tfrac12+\frac{s(s-1)}2
\int_1^\infty(x^{s/2}+x^{(1-s)/2})\psi(x)\,\frac{dx}{x}.
\]
The integral and all its $s$ derivatives converge locally
uniformly on the entire plane by exponential decay, so
analytic continuation proves this displayed identity for all $s$.
On $0\le\Re s\le1$, $|\Im s|\le3$, the product
$|s(s-1)|\le10$, the powers have combined modulus at most
$2x^{1/2}$, and $n^2\ge1+3(n-1)$ for positive integers.
Using $x^{-1/2}\le1$ for $x\ge1$ gives the full integral bound
\begin{equation}
\left|\xi(s)-\tfrac12\right|
\le\frac{10e^{-\pi}}{\pi(1-e^{-3\pi})}
<\frac{8000}{47994}<\frac13,\qquad \Re\xi(s)>\frac16.
\tag{AX5}
\end{equation}
The strict rational bound uses only $\pi>3$, $e^3>20$ and
$e^9>8000$. No tail or Gamma factor was discarded.
The classification AX4, the boundary argument and AX5 prove
\begin{equation}
\xi(r)\ne0\text{ at every finite root of }h_\varepsilon
\quad(\varepsilon>0),\qquad
j_{h_\varepsilon}\xi\text{ is a unit in }\C[s]/(h_\varepsilon).
\tag{AX6}
\end{equation}
AR7 retains root multiplicity at the double root. At
$\varepsilon^2=2$ this algebra has degree three; the infinity
residue is AR8, not a fictitious value $\xi(\infty)$.

\section{The precise obstruction to extending the selected residue}
The global infinity residue AR8 is not the continuous extension
of the selected finite-root residue AR3. This failure already
occurs at the original finite source point $q_{\mathrm M}$,
not only at the missing sign chart. For complex $a\ne0$ take
$q_a=(a,0,i/2,0)$, so $q_a\to q_{\mathrm M}=(0,0,i/2,0)$
and the selected root is $-i/a$. For real $\varepsilon>0$,
the completion, reflection and the exact Gamma identity
$|\Gamma(1/2+it)|^2=\pi/\cosh(\pi t)$ give the equality
below for all $\varepsilon>0$, and its upper bound for
$0<\varepsilon\le1$:
\begin{equation}
\left|\varepsilon^2\xi(-i/\varepsilon)\right|
=\frac{\sqrt{1+\varepsilon^2}}
       {2\sqrt{\cosh(\pi/(2\varepsilon))}}
 |\zeta(1+i/\varepsilon)|
\le\sqrt{\frac{1+\varepsilon^2}{2}}\,
 [\log(1+\varepsilon^{-1})+4]e^{-\pi/(4\varepsilon)}
\longrightarrow0 .
\tag{AX7}
\end{equation}
For completeness, the elementary bound used here follows by
summation against the integer counting function: for $\Re s>0$,
$s\ne1$, and integer $N\ge1$,
\[
\zeta(s)=\sum_{n=1}^{N}n^{-s}
       +\frac{N^{1-s}}{s-1}
       -s\int_N^\infty\{x\}x^{-s-1}\,dx .
\]
It follows first for $\Re s>1$ by integrating the tail counting
function, and then throughout the stated half-plane by uniform
convergence of the integral on compacta. At $s=1+iT$, $T\ge1$,
$N=\lceil T\rceil$, its three absolute bounds are
$1+\log N$, $1/T$, and $|1+iT|/N$, whose sum is less than
$\log(T+1)+4$. Since $\cosh t\ge e^t/2$, AX7 follows.
The Gamma identity follows directly from its reflection
formula at $1/2+it$, retaining the conjugate factor.

In contrast, take $a_n=-i/(2n)$, with $n$ a positive integer.
The same original source path has selected root $2n$ and
exact selected residue
\begin{equation}
a_n^2\xi(2n)
=-\frac{2n-1}{4n}\,(n-1)!\,\pi^{-n}\zeta(2n),
\qquad |a_n^2\xi(2n)|\longrightarrow\infty.
\tag{AX8}
\end{equation}
This follows by inserting $\Gamma(n)=(n-1)!$ in the original
completion. The lower bound $\zeta(2n)>1$ and the ratio
$n/\pi$ of successive factorial-over-$\pi$ factors prove
divergence. Both $q_{a_n}$ and $q_{1/(2n)}$ tend to the
same original point $q_{\mathrm M}$, but their selected
residues have incompatible limits. Hence AR3 has no continuous
extension there. The same conclusion holds at the negative
factor state in the full completion: its residue is identical
for each $a$, and its six factors have the opposite finite limit.
At the limiting target $h_0=s^3-s$, AR9 instead gives the
global residue $\Res_\infty=(3-\pi)/12\ne0$.
It therefore cannot repair this noncontinuity by being assigned
as the missing selected-root value. The exact relation that
does hold is the global residue identity AR8 and the complete
finite remainder map AR6, with all marked signs retained.
Every fixed injective receiving map transports these two
sequences and their unequal arithmetic limits unchanged;
no change of the original Gamma norm removes the obstruction.

The exact certificate
\texttt{verify\_arithmetic\_residues\_and\_jets.py}
checks the original factors, both residue charts, the quartet
matrix and inverse, the confluent evaluation and Hermite inverse,
the principal-part constants and the original escape polynomial.
It treats transcendental values and derivatives as formal scalars.

\begin{thebibliography}{9}
\bibitem{ES488} The Clankers, \emph{Erd\H{o}s--Straus project reader},
canonical 488-page source version 69, theorem
\texttt{explicit-four-time-Jacobian-collision}, lines 19484--19588.
Original-source SHA-256:
\texttt{678b3f70ee7d2a6d00f6c588d7778e390118ab4743a5759612dbba6430984c78}.
This is the provenance of the original polynomial and four
marked states, not attribution of the new arithmetic calculation.
\bibitem{Global} Programme calculation,
\emph{Complete fibres, the eight-sheet domain, and the nonproper-value
locus of the ES--Fable map},
\texttt{GLOBAL\_FIBRE\_AND\_NONPROPER\_LOCUS.tex}, GF1--9 and GF19--28.
Its factor inverse and infinity signs are rederived above.
\bibitem{Incoming} Incoming AI-generated web-session calculation,
equations (1)--(17) and the unnumbered zeta determinant following
(16), received 20 September 2026, attachment
\texttt{a422456a-2027-4452-bcfe-e55baf44bde0/Pasted text.txt}.
The incoming arithmetic remainders and inverse are independently
proved here; AR13 corrects its zeta determinant.
\bibitem{DLMFTheta} W. P. Reinhardt and P. L. Walker,
\emph{Theta Functions}, NIST Digital Library of Mathematical Functions,
\href{https://dlmf.nist.gov/20.2.E3}{20.2.3} and
\href{https://dlmf.nist.gov/20.7.E32}{20.7.32}.
Unchanged original equation TeX files
\texttt{20.2.E3.tex}, \texttt{20.7.E32.tex}
are retained in the parent \texttt{residue\_intake\_20260920} source set.
Their actual series and modular transformation derive the full
theta integral used in AX5.
\bibitem{DLMFXi} T. M. Apostol,
\emph{Zeta and Related Functions}, NIST Digital Library of
Mathematical Functions,
\href{https://dlmf.nist.gov/25.4.E3}{25.4.3} and
\href{https://dlmf.nist.gov/25.4.E4}{25.4.4}.
Original equation TeX retained as \texttt{25.4.E3.tex} and
\texttt{25.4.E4.tex}. These give the reflection and exact
completion factors used throughout; DLMF's source notes refer
to Apostol (1976), p.~260. The classical line-one argument
associated with Hadamard and de la Vallée Poussin (1896) is
reproduced in full above rather than attributed to this calculation.
\bibitem{DLMFGamma} R. A. Askey and R. Roy,
\emph{Gamma Function}, NIST Digital Library of Mathematical Functions,
\href{https://dlmf.nist.gov/5.4.E4}{5.4.4},
original equation TeX \texttt{5.4.E4.tex}; exact Gamma modulus
identity used in AX7, also derived there from reflection.
\end{thebibliography}
\end{document}
